% Part: first-order-logic
% Chapter: natural-deduction
% Section: identity

\documentclass[../../../include/open-logic-section]{subfiles}

\begin{document}

\olfileid[hi]{fol}{ntd}{ide}

\olsection{\usetoken{P}{derivation}: \usetoken{S}{identity} सहित}

!!^{derivation}s, जिनमें !!{identity} हो, अतिरिक्त अनुमान-नियम मांगती हैं।

\begin{defish}
\AxiomC{}
\RightLabel{\Intro{\eq}}
\UnaryInfC{$\eq[t][t]$}
\DisplayProof
\hfill
\begin{tabular}{r}
\AxiomC{$\eq[t_1][t_2]$}
\AxiomC{$!A(t_1)$}
\RightLabel{\Elim{\eq}}
\BinaryInfC{$!A(t_2)$}
\DisplayProof
\\[3ex]
\AxiomC{$\eq[t_1][t_2]$}
\AxiomC{$!A(t_2)$}
\RightLabel{\Elim{\eq}}
\BinaryInfC{$!A(t_1)$}
\DisplayProof
\end{tabular}
\end{defish}

उपर्युक्त नियमों में $t$, $t_1$, और $t_2$ बंद पद हैं। \Intro{\eq} नियम
$\eq[t][t]$ रूप के किसी भी तादात्म्य-कथन को बिना किसी मान्यता के सीधे
!!{derive} करने देता है।

\begin{ex}
यदि $s$ और $t$ बंद पद हैं, तो $!A(s), \eq[s][t] \Proves !A(t)$:
\begin{prooftree}
\AxiomC{$\eq[s][t]$}
\AxiomC{$!A(s)$}
\RightLabel{$\Elim{\eq}$}
\BinaryInfC{$!A(t)$}
\end{prooftree}
इसे ``अभिन्न वस्तुओं की प्रतिस्थापनीयता का सिद्धांत'' अथवा ``लाइबनित्स का
नियम'' भी कहते हैं।
\end{ex}

\begin{prob}
सिद्ध कीजिए कि $=$ सममित और संक्रामी दोनों है; अर्थात्
$\lforall[x][\lforall[y][(\eq[x][y] \lif
    \eq[y][x])]]$ और $\lforall[x][\lforall[y][\lforall[z]((\eq[x][y]
    \land \eq[y][z]) \lif \eq[x][z])]]$
की !!{derivation}s दीजिए।
\end{prob}

\begin{ex}
!!{derive} किया जाने वाला !!{sentence} है
\begin{align*}
& \lforall[x][\lforall[y][((!A(x) \land !A(y)) \lif \eq[x][y])]]
\intertext{और जिस !!{sentence} से इसे व्युत्पन्न करते हैं, वह है:}
& \lexists[x][\lforall[y][(!A(y) \lif \eq[y][x])]]
\end{align*}
हम !!{derivation} को पीछे की ओर निर्मित करते हैं:
\begin{prooftree}
\AxiomC{$\lexists[x][\lforall[y][(!A(y) \lif \eq[y][x])]]
\quad \Discharge{!A(a) \land !A(b)}{1}$}
\DeduceC{$\eq[a][b]$}
\DischargeRule{\Intro{\lif}}{1}
\UnaryInfC{$((!A(a) \land !A(b)) \lif \eq[a][b])$}
\RightLabel{\Intro{\lforall}}
\UnaryInfC{$\lforall[y][((!A(a) \land !A(y)) \lif \eq[a][y])]$}
\RightLabel{\Intro{\lforall}}
\UnaryInfC{$\lforall[x][\lforall[y][((!A(x) \land !A(y)) \lif \eq[x][y])]]$}
\end{prooftree}
अब मुख्य मान्यता का उपयोग करना होगा: वह एक अस्तित्वपरक !!{formula} है, इसलिए
मध्यवर्ती निष्कर्ष $\eq[a][b]$ को !!{derive} करने के लिए हम
\Elim{\lexists} का उपयोग करते हैं।
\begin{prooftree}
\AxiomC{$\lexists[x][\lforall[y][(!A(y) \lif \eq[y][x])]]$}
\AxiomC{$\Discharge{\lforall[y][(!A(y) \lif \eq[y][c]])}{2}$}
\noLine
\UnaryInfC{$\Discharge{!A(a) \land !A(b)}{1}$}
\DeduceC{$\eq[a][b]$}
\DischargeRule{\Elim{\lexists}}{2}
\BinaryInfC{$\eq[a][b]$}
\DischargeRule{\Intro{\lif}}{1}
\UnaryInfC{$((!A(a) \land !A(b)) \lif \eq[a][b])$}
\RightLabel{\Intro{\lforall}}
\UnaryInfC{$\lforall[y][((!A(a) \land !A(y)) \lif \eq[a][y])]$}
\RightLabel{\Intro{\lforall}}
\UnaryInfC{$\lforall[x][\lforall[y][((!A(x) \land !A(y)) \lif \eq[x][y])]]$}
\end{prooftree}
ऊपर दाईं ओर की उप-!!{derivation} अपनी मान्यताओं से $\eq[a][c]$ और
$\eq[b][c]$ दिखाकर पूरी होती है। इसके लिए दो अलग-अलग !!{derivation}s चाहिए।
$\eq[a][c]$ की !!{derivation} इस प्रकार है:
\begin{prooftree}
\AxiomC{$\Discharge{\lforall[y][(!A(y) \lif \eq[y][c]])}{2}$}
\RightLabel{\Elim{\lforall}}
\UnaryInfC{$!A(a) \lif \eq[a][c]$}
\AxiomC{$\Discharge{!A(a) \land !A(b)}{1}$}
\RightLabel{\Elim{\land}}
\UnaryInfC{$!A(a)$}
\RightLabel{\Elim{\lif}}
\BinaryInfC{$\eq[a][c]$}
\end{prooftree}
$\eq[a][c]$ और $\eq[b][c]$ से हम \Elim{\eq} द्वारा $\eq[a][b]$
!!{derive} करते हैं।
\end{ex}

\begin{prob}
ऐसी !!{derivation}s दीजिए कि इनमें से प्रत्येक का निष्कर्ष निम्नलिखित
!!{formula}s में से एक हो:
\begin{enumerate}
\item $\lforall[x][\lforall[y][((\eq[x][y] \land !A(x)) \lif !A(y))]]$
\item $\lexists[x][!A(x)] \land \lforall[y][\lforall[z][((!A(y) \land
    !A(z)) \lif \eq[y][z])]] \lif \lexists[x][(!A(x) \land
  \lforall[y][(!A(y) \lif \eq[y][x])])]$
\end{enumerate}
\end{prob}

\end{document}
